% Section 4.10, from lectures/L04/README.md.

\section{Summary}
\label{c4:sec:review}

\needspace{6\baselineskip}
\subsection*{What you should be able to do now}
\begin{itemize}
  \item Explain what metastability is, why an asynchronous input can cause it, and why two
    flip-flops make it unlikely rather than impossible.
  \item Apply the two-flip-flop synchroniser to any asynchronous input, including the ``assert
    asynchronously, release synchronously'' pattern used for reset.
  \item Say why that chain also works as a practical but incomplete debounce, and what a real one
    would add.
  \item Give a module a \code{generic}, override it with a \code{generic map}, write a port's width
    in terms of it, and judge when a generic earns its place instead of being added on principle.
\end{itemize}

\needspace{6\baselineskip}
\subsection*{Questions to test yourself with}
\begin{enumerate}
  \item What exactly happens when an asynchronous input to a flip-flop changes too close to an
    active clock edge?
  \item How does the two-flip-flop synchroniser handle that, and why is it the \emph{second}
    flip-flop's output that is safe to use?
  \item Why does the chain used here behave as a debounce in CircuitVerse but not fully debounce a
    real button at \code{50 MHz}, and why is that a property of the clock period rather than of the
    circuit? What would you add to fix it?
  \item \code{button_sync} has a generic and \code{reset_sync} has none. What is the rule, and what
    would be wrong with giving \code{reset_sync} a \code{WIDTH} generic ``for symmetry''?
\end{enumerate}

\needspace{6\baselineskip}
\subsection*{Further reading}
Erik Pihl's video walkthrough of metastability and the two-flip-flop synchroniser is at
\url{https://www.youtube.com/watch?v=KrssJRgF13I} as additional, optional viewing.
\url{https://nandland.com/lesson-13-metastability/} and
\url{https://vhdlwhiz.com/terminology/metastability/} cover the same concepts from a slightly
different angle, useful if this chapter does not land on the first reading.

\needspace{6\baselineskip}
\subsection*{Looking ahead}
\Chapref{c5:ch} goes on with:
\begin{itemize}
  \item \code{variable}: the one construct you have not yet met, and the sharp difference between
    how it updates and how a \code{signal} does.
  \item Why the same accumulation written with a \code{signal} and with a \code{variable} gives two
    different answers, and which of the two is the bug.
  \item What the toolchain actually builds: LUTs and flip-flops, what limits a clock's speed, and
    the latch you describe by accident.
\end{itemize}
